\frontchapter{II. The Book's Recurring Approach}
\section*{1. A Fixed Structure for Each Chapter}
\begin{enumerate}
\item \textbf{Pick up an object}: A stone tool, a coin, a diary page, a song, a ticket, or a phone.
\item \textbf{Enter a scene}: First place the reader in a market, temple, courtroom, port, theater, factory, or trench.
\item \textbf{Ask a real question}: Let a conflict emerge before offering a conclusion.
\item \textbf{Hear multiple voices}: Rulers and ordinary people, victors and the defeated, adults and children, centers and margins.
\item \textbf{Learn a tool}: A timeline, map, close reading, conceptual distinction, argument diagram, or visual analysis.
\item \textbf{Read a short piece of evidence}: A primary document, object, image, statistical record, or oral memory.
\item \textbf{Compare explanations}: Why the same fact can have different explanations that are not equally convincing.
\item \textbf{Go deeper}: Introduce theories from linguistics, narratology, religious studies, historiography, philosophy, or aesthetics.
\item \textbf{Return to today}: Consider how the old question persists in schools, families, cities, and online life.
\item \textbf{Leave an open question}: Require reasons, not a guess at the author's position.
\end{enumerate}
\section*{2. Dig Six Layers into One Question}
Take ``why can a joke hurt someone?'' as an example:

\begin{enumerate}
\item The same words have different effects when spoken by different people.
\item Words carry tone, relationships, and context as well as literal meaning.
\item Speaking itself can promise, command, exclude, or establish an identity.
\item Power within a group determines who can call something a joke and who must bear its consequences.
\item Literature uses irony and unreliable narrators to give a sentence several layers of meaning at once.
\item Philosophy asks further: does harm arise from the speaker's intention, actual consequences, or social rules?
\end{enumerate}
Accessibility does not mean staying at the first layer. It means always providing a step down to the next.

\section*{3. Three Reading Levels}
\begin{itemize}
\item \textbf{Main Story}: Driven by people, scenes, objects, and questions; middle school students can read it independently.
\item \textbf{Thinking Bridge}: Learn close reading, argument, source evaluation, maps, and simple statistics.
\item \textbf{Further Challenge}: Explore linguistics, narratology, hermeneutics, philosophy of religion, ethics, political philosophy, historiography, and aesthetics.
\end{itemize}
\section*{4. The Density of Examples}
Each chapter should generally include:

\begin{itemize}
\item 1 vivid historical or everyday scene;
\item 3--5 cross-cultural examples;
\item 1 short primary source;
\item 1 easily misjudged counterexample;
\item 1 exercise in stating the other side's reasons fully;
\item 1 map, timeline, concept map, or argument diagram;
\item 1 echo in contemporary life;
\item 1 doorway to deeper theory.
\end{itemize}
The book is expected to use 500--700 short examples. A few old friends return: a coin, a map, a trade route, a temple, a poem, a cup of coffee, a passport, and a phone. Each encounter uses a different lens.

\section*{5. Four Kinds of Content Must Be Distinguished}
\begin{itemize}
\item \textbf{What happened}: The range of facts supported by current evidence.
\item \textbf{Why it happened}: Causal explanations that may compete with one another.
\item \textbf{How people at the time understood it}: The conceptual world of historical actors.
\item \textbf{How we judge it}: The value judgments of readers and authors.
\end{itemize}
Explaining why an institution arose is not defending it. Understanding a belief does not require readers to share it. Criticizing an idea cannot be done by distorting it.

\section*{6. Principles of Global History and Multiple Perspectives}
\begin{itemize}
\item Never present a civilization as a sealed box that had no contact with others.
\item Include China, South Asia, West Asia, Africa, Europe, the Americas, and Oceania in the main narrative, not an ``other regions'' appendix.
\item Do not rank societies on a single scale of civilizational advancement.
\item Attend to emperors, merchants, farmers, craftspeople, women, children, enslaved people, migrants, and minorities together.
\item Compare similar questions without forcing different traditions to answer one specifically Western or Chinese question.
\item When using BCE/CE and regional dating systems, explain their different calendars and views of time.
\end{itemize}
